\documentclass
  [hyperref={colorlinks = true,linkcolor = blue, 
             citecolor = blue, urlcolor = blue}
  ]{beamer}

\setbeamertemplate{navigation symbols}{}
\setbeamertemplate{footline}[frame number]
 
\usepackage[utf8]{inputenc}

\usepackage[newfloat]{minted}

\usepackage{pgf}
\usepackage{tikz}
\usepackage{upquote}
\usepackage{natbib}
\usepackage[export]{adjustbox}
\usepackage{amsmath}
\usepackage{amssymb}
\usepackage{trfrac}

\usetikzlibrary{arrows,automata,fit, shapes.geometric}

\bibliographystyle{abbrvnat}

\newenvironment{code}{\captionsetup{type=listing}}{}
\SetupFloatingEnvironment{listing}{name=Listing}

\setbeamertemplate{items}[square]

\title{Towards Palatable Functional Programming with 
       Dependent Type Theories}
\author{Donovan Crichton}
  
\date{August 2023}

\begin{document}
 
\frame{\titlepage}

\begin{frame}[fragile]
  \frametitle{Preliminaries}
  \begin{itemize}
  \item Slides and Examples available at:
    \href{https://github.com/donovancrichton/Talks}
         {https://github.com/donovancrichton/Talks}
  \item This talk: BFPG/PalatableProgrammingWithDTT
  \end{itemize}
\end{frame}

\begin{frame}[fragile]
\frametitle{About me}
\begin{minipage}[t]{\linewidth}
\noindent
  \begin{minipage}{0.5\linewidth}
    \begin{figure}[h]
       \includegraphics[width=\linewidth]
         {Images/ANU.png}
    \end{figure}
  \end{minipage}
  \hfill
  \begin{minipage}{0.45\linewidth}
    \begin{itemize}
      \item 
        \href{https://cecc.anu.edu.au/people/donovan-crichton}
             {PhD Candidate}
      \item 
        \href{https://comp.anu.edu.au/research/clusters/computing-foundations/}
             {Computing Foundations}
      \item 
        \href{https://comp.anu.edu.au/}
             {School of Computing}
    \end{itemize}
  \end{minipage}
\end{minipage}
\hrule
  \begin{minipage}{0.5\linewidth}
    \begin{figure}[h]
      \includegraphics[scale=0.5,width=\linewidth]
        {Images/GriffithLogo.png}
    \end{figure}
  \end{minipage}
  \hfill
  \begin{minipage}{0.45\linewidth}
  \begin{itemize}
    \item Visiting Scholar
    \item Trusted Systems Lab
    \item \href{https://www.griffith.edu.au/institute-integrated-intelligent-systems}{IIIS}
  \end{itemize}
  \end{minipage}
\hrule
  \begin{minipage}{0.5\linewidth}
    \begin{figure}[h]
      \includegraphics[width=\linewidth]
        {Images/asd-logo.png}
    \end{figure}
  \end{minipage}
  \hfill
  \begin{minipage}{0.45\linewidth}
  \begin{itemize}
    \item ASD 
      \href{https://www.asd.gov.au/about/asd-anu-co-lab}{Co-Lab} 
      Scholar
  \end{itemize}
  \end{minipage}
\end{frame}

\begin{frame}[fragile]
  \frametitle{Rose Tree Maps: Haskell vs Idris}
  The Idris code will not pass the totality checker.
  \begin{minipage}{0.44\linewidth}
    \inputminted{haskell}{../Code/src/HaskellRose.hs}
  \end{minipage}
  $\;$
  \vline
  $\;$
  \begin{minipage}{0.5\linewidth}
    \inputminted{idris}{../Code/src/IdrisRoseBad.idr}
  \end{minipage}
\end{frame}

\begin{frame}[fragile]
\frametitle{What is a dependent type?}
\setbeamercolor{block title}{fg=white,bg=orange!75!black}
\setbeamercolor{block body}{fg=white,bg=orange!25!black}
\begin{block}{Dependent Type(s): $\Pi_{x \in A}.B(x)$ instead of 
 $\lambda_{x \in A}.B$}
A dependent type is a type that \textit{depends} on a specific
value of its input.
\end{block}
\setbeamercolor{block title}{fg=white,bg=purple!75!black}
\setbeamercolor{block body}{fg=black,bg=white!80!gray}
\begin{block}{$\lambda_{x \in A}.B$ in Idris code, in general.}
\begin{minted}{idris}
f : a -> b
f x = ?somedefinition
\end{minted}
\end{block}
\setbeamercolor{block title}{fg=white,bg=purple!75!black}
\setbeamercolor{block body}{fg=black,bg=white!80!gray}
\begin{block}{$\Pi_{x \in A}.B(x)$ in Idris code, in general.}
\inputminted{idris}{../Code/src/PiGeneralExample.idr}
\end{block}
\end{frame}

\begin{frame}[fragile]
\frametitle{What is a dependent type? (concrete examples)}
\setbeamercolor{block title}{fg=white,bg=purple!75!black}
\setbeamercolor{block body}{fg=black,bg=white!80!gray}
\begin{block}{$\lambda_{x \in A}.B$ in Idris code, concretely.}
\begin{minted}{idris}
f : a -> List 
f x = [x]
\end{minted}
\end{block}
\setbeamercolor{block title}{fg=white,bg=purple!75!black}
\setbeamercolor{block body}{fg=black,bg=white!80!gray}
\begin{block}{$\Pi_{x \in A}.B(x)$ in Idris code, concretely.}
  \inputminted{idris}{../Code/src/PiSpecificExample.idr}
\end{block}

\end{frame}

\begin{frame}[fragile]
\frametitle{What is a dependent type? (diagrammatically)}
\begin{tikzpicture}
% Boolean set
\node (BTrue) at (0,7) {True};
\node[below of=BTrue] (BFalse) {False};
% Character set
\node (StrA) at (6, 9) {`a'};
\node[right of=StrA] (StrB) {`b'};
\node[right of=StrB] (StrC) {`c'};
\node[right of=StrC] (StrRest) {\ldots};
% Natural Number set
\node (Num1) at (6, 4) {0};
\node[right of=Num1] (Num2) {1};
\node[right of=Num2] (Num3) {2};
\node[right of=Num3] (NumRest) {\ldots};
% arrows
\path[->] (BTrue) edge (StrB);
\path[->] (BFalse) edge (Num2);
% set ellipses
\node[shape=ellipse,draw,minimum size=2cm,
      fit={(StrA) (StrRest)}] {};
\node[shape=ellipse,draw,minimum size=2cm,
      fit={(Num1) (NumRest)}] {};
\node[shape=ellipse,draw,minimum size=2cm,
      fit={(BTrue) (BFalse)}] {};
\end{tikzpicture}
\end{frame}

\begin{frame}[fragile]
\frametitle{What is a Dependent Type? (diagrammatic example)}

\setbeamercolor{block title}{fg=white,bg=purple!75!black}
\setbeamercolor{block body}{fg=black,bg=white!80!gray}
\begin{block}{Did you guess the right Idris type?}
  \inputminted{idris}{../Code/src/PiDiagramExample.idr}
\end{block}

\setbeamercolor{block title}{fg=white,bg=green!55!black}
\setbeamercolor{block body}{fg=white,bg=green!25!black}
\begin{block}{What about $\Sigma_{x \in A}.B(x)$?}
Dependent pairs can be constructed from $\Pi$ types if
your language supports recursive algebraic data type
definitions.
\end{block}
\end{frame}

\begin{frame}[fragile]
\frametitle{Why do we mean by totality? (good)}
\setbeamercolor{block title}{fg=white,bg=orange!75!black}
\setbeamercolor{block body}{fg=white,bg=orange!25!black}
\begin{block}{Total Functions}
A functions is \textit{total} if it terminates and is defined
for all possible values for its domain.
\end{block}
\setbeamercolor{block title}{fg=white,bg=purple!75!black}
\setbeamercolor{block body}{fg=black,bg=white!80!gray}
\begin{block}{Functions that are total.}
  \inputminted{idris}{../Code/src/TotalExample.idr}
\end{block}
\end{frame}

\begin{frame}[fragile]
\frametitle{Why do we mean by totality? (bad)}
\setbeamercolor{block title}{fg=white,bg=orange!75!black}
\setbeamercolor{block body}{fg=white,bg=orange!25!black}
\begin{block}{Total Functions}
A functions is \textit{total} if it terminates and is defined
for all possible values for its domain.
\end{block}
\setbeamercolor{block title}{fg=white,bg=purple!75!black}
\setbeamercolor{block body}{fg=black,bg=white!80!gray}
\begin{block}{Functions that are \textit{not} total.}
  \inputminted{idris}{../Code/src/TotalBadExample.idr}
\end{block}
\end{frame}

\begin{frame}[fragile]
\frametitle{Why do we care about totality?}
\setbeamercolor{block title}{fg=white,bg=green!55!black}
\setbeamercolor{block body}{fg=white,bg=green!25!black}
\begin{block}{Can we type check this?}
$f : \mathbb{N} \rightarrow \mathbb{N}$ \\
$f = \lambda(x \in \mathbb{N}).f\;x$ \\
$\;$\\
$v : Vect\;k\;\alpha -> Vect\; (f\;k)\;\alpha$ \\
$v\;xs = xs$
\end{block}

\setbeamercolor{block title}{fg=white,bg=purple!75!black}
\setbeamercolor{block body}{fg=black,bg=white!80!gray}
\begin{block}{Infinite recursion in types in Idris}
  \inputminted{idris}{../Code/src/ForeverExample.idr}
\end{block}
\end{frame}

\begin{frame}[fragile]
\frametitle{Using the totality pragma.}

\setbeamercolor{block title}{fg=white,bg=orange!75!black}
\setbeamercolor{block body}{fg=white,bg=orange!25!black}
\begin{block}{How Idris reduces $\Pi$ application in types.}
Idris will \textit{only} reduce function application in types
if the function has passed the totality checker.
\end{block}

\setbeamercolor{block title}{fg=white,bg=purple!75!black}
\setbeamercolor{block body}{fg=black,bg=white!80!gray}
\begin{block}{The totality checker blocks evaluation:}
  \inputminted{idris}{../Code/src/ForeverExampleFail.idr}
\end{block}
\end{frame}

\begin{frame}[fragile]
\frametitle{Back to our rose tree example.}
\setbeamercolor{block title}{fg=white,bg=green!55!black}
\setbeamercolor{block body}{fg=white,bg=green!25!black}
\begin{block}{Why does our rose tree Functor instance fail?}
If map on lists is total, and our RoseTree is total. Why do we get
errors when writing the functor instance?
\end{block}

\setbeamercolor{block title}{fg=white,bg=purple!75!black}
\setbeamercolor{block body}{fg=black,bg=white!80!gray}
\begin{block}{Map on lists is total:}
  \inputminted{idris}{../Code/src/TotalMapList.idr}
\end{block}
\end{frame}

\begin{frame}[fragile]
\frametitle{Back to our rose tree example (2).}
\setbeamercolor{block title}{fg=white,bg=purple!75!black}
\setbeamercolor{block body}{fg=black,bg=white!80!gray}
\begin{block}{The definition of rose tree is total}
  \inputminted{idris}{../Code/src/TotalMapList.idr}
\end{block}
\setbeamercolor{block title}{fg=white,bg=green!55!black}
\setbeamercolor{block body}{fg=white,bg=green!25!black}
\begin{block}{Is there a problem with our map definition?}
\texttt{map f (MkRose x xs) = MkRose (f x) (map (map f) xs)}
\end{block}
\end{frame}

\begin{frame}[fragile]
  \frametitle{How totality checking actually works}
  \setbeamercolor{block title}{fg=white,bg=orange!75!black}
  \setbeamercolor{block body}{fg=white,bg=orange!25!black}
  \begin{block}{Syntactic Guardedness}
  Idris (and all other dependently typed proof assistants) check
  termination by means of a \textit{syntactic guard} condition.
  \end{block}
  \setbeamercolor{block title}{fg=white,bg=purple!75!black}
  \setbeamercolor{block body}{fg=black,bg=white!80!gray}
  \begin{block}{Recursive calls must be \textit{guarded} by a
  destructor.}
    \inputminted{idris}{../Code/src/GuardedRecursiveCall.idr}
  \end{block}
\end{frame}

\begin{frame}[fragile]
  \frametitle{How totality checking actually works (2)}
  \setbeamercolor{block title}{fg=white,bg=purple!75!black}
  \setbeamercolor{block body}{fg=black,bg=white!80!gray}
  \begin{block}{Corecursive calls must be \textit{guarded} by a
  constructor, and correctly \textit{annotated} as Inf.}
    \inputminted{idris}{../Code/src/GuardedCorecursiveCall.idr}
  \end{block}
\end{frame}

\begin{frame}[fragile]
  \frametitle{Summary of totality via syntactic guards}
  \setbeamercolor{block title}{fg=white,bg=orange!55!black}
  \setbeamercolor{block body}{fg=white,bg=orange!25!black}
  \begin{block}{The \texttt{(map f)} call on our RHS does not
    meet the syntactic guard restriction.}
  \texttt{map f (MkRose x xs) = MkRose (f x) (map (map f) xs)}
  \end{block}
  \begin{itemize}
    \item Syntactic guards do not work well with HOF.
    \item Inf for corecursion is broken.
    \item Most corecursive approaches are broken (sized types,
    delay/force, infinity) what about co-patterns?
    \item Manual proofs of well-foundedness can be onerous.
  \end{itemize}
\end{frame}

\begin{frame}[fragile]
  \frametitle{Introducing Type Theory}
  \setbeamercolor{block title}{fg=white,bg=orange!75!black}
  \setbeamercolor{block body}{fg=white,bg=orange!25!black}
  \begin{block}{Components of a Type Theory}
  A Type Theory consists of a grammar, a set of judgement rules, 
  a set of desirable properties that are proven to hold, and
  a given semantics.
  \end{block}
  \setbeamercolor{block title}{fg=white,bg=green!75!black}
  \setbeamercolor{block body}{fg=white,bg=green!25!black}
  \begin{block}{Why do we need all this?}
  In a dependently typed setting, the PAT (Propositions-As-Types)
  interpretation is assumed to hold. Alterations to any of these
  features (i.e extra syntactic annotations, different type
  checking algorithms, etc) can break the desirable properties
  and in the worst case, allow you to prove things that are not
  true.
  \end{block}

    %grammar
  %rules
  %proofs of desirable properties
  %consistency
  %canonicity
  %normalization
  %decidability
  %attached semantics
\end{frame}

\begin{frame}[fragile]
  \frametitle{Grammar (Parsed Lexemes)}
  \setbeamercolor{block title}{fg=white,bg=orange!75!black}
  \setbeamercolor{block body}{fg=white,bg=orange!25!black}
  \begin{block}{What is a grammar?}
    A formal set of syntax used in the definition of your
    type theory, usually given in BNF.
  \end{block}
  \setbeamercolor{block title}{fg=white,bg=purple!75!black}
  \setbeamercolor{block body}{fg=black,bg=white!80!gray}
  \begin{block}{An example grammar for a dependent type theory}
  \begin{align*}
    \Gamma ::=&\;\diamond \;|\; \Gamma, x : S \\
    S, T, U ::=&\;(x : S) \rightarrow T \;|\; (x : S) \otimes T 
                \;|\; El(M) \\& \;|\; Bool \;|\; Unit 
                  \;|\; \mathcal{U} \\
    P, M, N ::= &\;x \;|\; \Pi (x : S).M \;|\; M \; N \;|\; True
                  \;|\; False \\
             &\;|\; \star \;|\; fst\;M \;|\; snd\;M 
                  \;|\; If\;P_{Bool} \; M \; N
  \end{align*}
  \end{block}
  
\end{frame}

\begin{frame}[fragile]
  \frametitle{Rules - Type Formers (Type Constructors)}
  \begin{minipage}{\linewidth}
  \setbeamercolor{block title}{fg=white,bg=orange!75!black}
  \setbeamercolor{block body}{fg=white,bg=orange!25!black}
  \begin{block}{Type Formers}
   A type formation rule is analogous to a type constructor.
  \end{block}
  \setbeamercolor{block title}{fg=white,bg=purple!75!black}
  \setbeamercolor{block body}{fg=black,bg=white!80!gray}
  \begin{block}{List Type Constructor in Idris}
  \begin{minted}{idris}
  data List : (s : Type) -> Type where
  \end{minted}
  \end{block}
  \end{minipage}
  \[
   \trfrac[\small(\emph{ListTy})]
     {\Gamma \vdash S : \mathcal{U}}
     {\Gamma \vdash \text{List } S : \mathcal{U}}
  \]
  
\end{frame}

\begin{frame}[fragile]
  \frametitle{Rules - Term Formers (Data Constructors)}
  \begin{minipage}{\linewidth}
  \setbeamercolor{block title}{fg=white,bg=orange!75!black}
  \setbeamercolor{block body}{fg=white,bg=orange!25!black}
  \begin{block}{Term Formers}
   A term formation rule is analogous to a data constructor.
  \end{block}
  \setbeamercolor{block title}{fg=white,bg=purple!75!black}
  \setbeamercolor{block body}{fg=black,bg=white!80!gray}
  \begin{block}{List Data Constructors in Idris}
  \begin{minted}{idris}
  data List : (s : Type) -> Type where
    Nil : List a
    (::) : a -> List a -> List a
  \end{minted}
  \end{block}
  \end{minipage}
  \[
   \trfrac[\small(\emph{Nil})]
     {\Gamma \vdash S : \mathcal{U} \qquad \Gamma 
      \vdash \text{List } S : \mathcal{U}}
     {\Gamma \vdash \left[ \right] : \text{List } S}
  \]
  \[
   \trfrac[\small(\emph{Cons})]
     {\Gamma \vdash S : \mathcal{U} \qquad 
      \Gamma \vdash \text{List } S : \mathcal{U}
      \qquad \Gamma \vdash x : S \qquad \Gamma \vdash
      xs : \text{List } S}
     {\Gamma \vdash x :: xs : \text{List } S}
  \]
\end{frame}

\begin{frame}[fragile]
  \frametitle{Properties}
  \setbeamercolor{block title}{fg=white,bg=orange!75!black}
  \setbeamercolor{block body}{fg=white,bg=orange!25!black}
  \begin{block}{Normalisation}
  Prove that all terms reduce to an irreducible normal form.
  I.e show that evaluation of terms eventually stops.
  \end{block}
  \setbeamercolor{block title}{fg=white,bg=orange!75!black}
  \setbeamercolor{block body}{fg=white,bg=orange!25!black}
  \begin{block}{Canonicity}
  Prove that every closed computation (evaluation in an empty
  context) of a type results in a singular canonical value.
  i.e for every term of $\diamond \vdash x : \mathbb{N}$ 
  we get a natural number literal when we evaluate.
  \end{block}
  \setbeamercolor{block title}{fg=white,bg=orange!75!black}
  \setbeamercolor{block body}{fg=white,bg=orange!25!black}
  \begin{block}{Consistency}
  Prove that we cannot derive $\bot$, which in dependent 
  type theories is the type with no terms.
  \end{block}
  \setbeamercolor{block title}{fg=white,bg=orange!75!black}
  \setbeamercolor{block body}{fg=white,bg=orange!25!black}
  \begin{block}{Decidability}
  Give a decision procedure for type checking, i.e an
  algorithm that for all $\Gamma \vdash x : A$ we can
  determine if $x \in A$ or not.
  \end{block}

\end{frame}

\begin{frame}[fragile]
  \frametitle{Semantics}
  \setbeamercolor{block title}{fg=white,bg=orange!75!black}
  \setbeamercolor{block body}{fg=white,bg=orange!25!black}
  \begin{block}{If we're just manipulating symbols, how can we
  prove theorems in `real' mathematics?}
  We need a \textit{semantics} for our syntax. A mapping of some
  area of mathematics to our symbols and back. Common semantics
  are denotational, operational, and axiomatic. See e.g 
  \citep{winskel1993formal}
  \end{block}
  \setbeamercolor{block title}{fg=white,bg=green!55!black}
  \setbeamercolor{block body}{fg=white,bg=green!25!black}
  \begin{block}{Which one do we use?}
  One of many category theoretic semantics used for dependent
  types is categories with families...also used are categories 
  with attributes, comprehension
  categories, display map categories, natural models, contextual
  categories. \\ Sadly beyond the scope of this talk (and me, for
  the moment!)
  \end{block}
\end{frame}

\begin{frame}[fragile]
  \frametitle{Guarded Dependent Type Theory}
  \setbeamercolor{block title}{fg=white,bg=green!55!black}
  \setbeamercolor{block body}{fg=white,bg=green!25!black}
  \begin{block}{All this talk of type theory...what about our
  Totality problem?}
  A promising alternative approach may be a type-directed
  approach to guardedness, instead of a syntactic one.
  \end{block}

  \setbeamercolor{block title}{fg=white,bg=orange!75!black}
  \setbeamercolor{block body}{fg=white,bg=orange!25!black}
  \begin{block}{Guarded Dependent Type 
    Theory \citep{bizjak2016guarded}}
  ``We present guarded dependent type theory, gDTT, an 
    extensional dependent type theory with a ‘later’ modality 
    and clock quantifiers for programming and proving with guarded
    recursive and coinductive types. The later modality is 
    used to ensure the productivity of recursive definitions in 
    a modular, type based, way.''
  \end{block}


\end{frame}

\begin{frame}[fragile]
  \frametitle{Extensionality and Intensionality}
  \setbeamercolor{block title}{fg=white,bg=orange!75!black}
  \setbeamercolor{block body}{fg=white,bg=orange!25!black}
  \begin{block}{Extensionality}
  Two objects may be judged to be equal if they have the same
  external properties.
  Best seen in functional extensionality: $f = g \iff \forall x, f(x) = g(x)$
  \end{block}
  
  \setbeamercolor{block title}{fg=white,bg=orange!75!black}
  \setbeamercolor{block body}{fg=white,bg=orange!25!black}
  \begin{block}{Intensionality}
  Two objects may be judged to be equal if they have the same
  internal properties. In practice, two expressions may be 
  judged equal if they reduce to the same normal form.
  \end{block}

  \setbeamercolor{block title}{fg=white,bg=purple!75!black}
  \setbeamercolor{block body}{fg=black,bg=white!80!gray}
  \begin{block}{Intensional Equality in Idris}
  \inputminted{idris}{../Code/src/IntensionalEquality.idr}
  \end{block}
\end{frame}

\begin{frame}[fragile]
  \frametitle{Observational Type Theory}
  \setbeamercolor{block title}{fg=white,bg=green!55!black}
  \setbeamercolor{block body}{fg=white,bg=green!25!black}
  \begin{block}{GDTT is extensional but Idris is
  intensional?}
  Fortunately, not all is lost. There have been recent advanced
  in other dependent type theories that promise \textit{just}
  enough extensionality to get the job done.
  \end{block}

  \setbeamercolor{block title}{fg=white,bg=orange!75!black}
  \setbeamercolor{block body}{fg=white,bg=orange!25!black}
  \begin{block}{Observational Type 
    Theory \citep{altenkirch2007observational} \\
    \citep{pujet2022observational}}
    ``Building on the recent extension of dependent type 
      theory with a universe of definitionally 
      proof-irrelevant types, we introduce $TT_{obs}$, a 
      new type theory based on the setoidal interpretation 
      of dependent type theory. $TT_{obs}$ equips every type 
      with an identity relation that satisfies 
      function extensionality, propositional extensionality, 
      and definitional uniqueness of identity proofs (UIP).''
  \end{block}
\end{frame}

\begin{frame}[fragile]
  \frametitle{Quantitative Type Theory}
  \setbeamercolor{block title}{fg=white,bg=green!55!black}
  \setbeamercolor{block body}{fg=white,bg=green!25!black}
  \begin{block}{Idris has quantities?}
  Idris is (mostly) based on a linear dependent type theory
  with the hope that programmers can use this to formally
  prove properties about execution times and memory usage.
  This means we're still missing a step!
  \end{block}

  \setbeamercolor{block title}{fg=white,bg=orange!75!black}
  \setbeamercolor{block body}{fg=white,bg=orange!25!black}
  \begin{block}{Quantitative Type Theory \citep{atkey2018syntax}}
  ``We present Quantitative Type Theory, a Type Theory that
  records usage information for each variable in a judgement,
  based on a previous system by McBride. The usage information is
  used to give a realizability semantics using a variant of
  Linear Combinatory Algebras, refining the usual realizability
  semantics of Type Theory by accurately tracking resource
  behaviour.''
  \end{block}
\end{frame}


\begin{frame}[fragile]
  \frametitle{Towards $GOLD_{tt}$ - Construction via Formalisation}
  \setbeamercolor{block title}{fg=white,bg=green!55!black}
  \setbeamercolor{block body}{fg=white,bg=green!25!black}
  \begin{block}{\textit{Finally! I can talk about my PhD!}}
  Work towards a Guarded, Observational, Linear, Dependent, Type
  Theory ($GOLD_{tt}$) that allows us to combine features from
  all theory type theories.
  \end{block}
  \setbeamercolor{block title}{fg=white,bg=orange!75!black}
  \setbeamercolor{block body}{fg=white,bg=orange!25!black}
  \begin{block}{Formalisation}
  To \textit{Formalise} or \textit{Mechanise} a type theory
  is to encode it in a proof assistant where machine checked
  proofs are given for as many of the components that make
  up a type theory as possible.
  \end{block}
  \setbeamercolor{block title}{fg=white,bg=green!55!black}
  \setbeamercolor{block body}{fg=white,bg=green!25!black}
  \begin{block}{Methodology}
  Formalise each type theory in turn to get deep insights about
  the inner workings. \\
  Then Formalise $GOLD_{tt}$. \\
  Then Give category theoretic semantics.
  \end{block}
\end{frame}

\begin{frame}[fragile]
  \frametitle{Closing and ANU Plug}
  \setbeamercolor{block title}{fg=white,bg=green!55!black}
  \setbeamercolor{block body}{fg=white,bg=green!25!black}
  \begin{block}{ANU is looking for PhD Students in Formal
  Methods!}
  The Computing Foundations cluster at the School of Computing, The
Australian National University, have multiple PhD scholarships available in
areas relevant to the types community, including logic, programming
languages, systems, formal methods, theory and software engineering.

If you want to do cutting edge research in any of these fields, or a mix
thereof, visit our
\href{https://comp.anu.edu.au/research/clusters/computing-foundations/}{page}
and get in touch with any of the staff listed. \\ 
Please also contact Michael
Norrish (cluster lead) for further details.
  \end{block}
Thank you!
\end{frame}

\begin{frame}[fragile]
\frametitle{References}
\bibliography{references}{}
\end{frame}




\end{document}

